\SourceLeaf{222}
\[
\begin{aligned}
p&=\tfrac{1}{m}(-\mathfrak{A}t+\mathfrak{B}u),&
q&=\tfrac{1}{m}(-\mathfrak{C}t+\mathfrak{D}u)\\
p&=-\tfrac{1}{m}(\mathfrak{A}t+\mathfrak{B}u),&
q&=-\tfrac{1}{m}(\mathfrak{C}t+\mathfrak{D}u)
\end{aligned}
\]
ita ut multitudo omnium formularum nunc quater maior sit quam antea, $t$ et $u$
vero non amplius omnes numeros aequationi $tt-(bb-ac)uu=mm$ satisfacientes exprimant,
sed positivos tantum. Quaevis harum formarum itaque seorsim considerari,
et qui valores ipsorum $t$, $u$ praebeant valores integros ipsorum $x$, $y$,
investigari debebit.

Ex formula
\[
p=\tfrac{1}{m}(\mathfrak{A}t+\mathfrak{B}u),\quad
q=\tfrac{1}{m}(\mathfrak{C}t+\mathfrak{D}u)\qquad .\ .\ .\ .\ [1]
\]
sequuntur valores ipsorum $x$, $y$ hi:
\[
x=\frac{\mathfrak{A}t+\mathfrak{B}u+mcd-mbe}{m(bb-ac)},\quad
y=\frac{\mathfrak{C}t+\mathfrak{D}u+mae-mbd}{m(bb-ac)}
\]
Supra vero ostendimus, omnes valores (positivos) ipsorum $t$ constituere progressionem
recurrentem $t^0,t',t''$ etc., similiter valores respondentes ipsius $u$ quoque
seriem recurrentem formare $u^0,u',u''$ etc.; praeterea assignari posse numerum $\rho$
talem, ut secundum modulum quemcunque datum fiat
\[
t^\rho\equiv t^0,\quad t^{\rho+1}\equiv t',\quad t^{\rho+2}\equiv t''\ \text{etc.},\quad
u^\rho\equiv u^0,\quad u^{\rho+1}\equiv u'\ \text{etc.}
\]
Pro hoc modulo accipiemus numerum $m(bb-ac)$, designabimusque brevitatis gratia
valores ipsorum $x$, $y$ qui prodeunt ponendo $t=t^0$, $u=u^0$, et quibus
tribuemus indicem $0$, per $x^0,y^0$; similiterque eos qui prodeunt faciendo
$t=t'$, $u=u'$, per $x',y'$ quibus tribuemus indicem $1$, etc.
Tunc nullo negotio perspicietur, si $x^h,y^h$ fuerint numeri integri atque $\rho$
rite determinatus, etiam $x^{h+\rho},y^{h+\rho}$; nec non
$x^{h+2\rho},y^{h+2\rho}$ et generaliter $x^{h+k\rho},y^{h+k\rho}$,
integros fore; et contra si $x^h$ vel $y^h$ sit fractus, etiam $x^{h+k\rho}$,
vel $y^{h+k\rho}$ fractum fore. Hinc facile concluditur, si valores ipsorum
$x$, $y$, quibus indices $0,1,2\,.\,.\,.\,.\,\rho-1$ competunt, evolvantur,
et pro nullo horum indicum \textit{tum} $x$, \textit{tum} $y$ integer sit,
nullum omnino indicem dari, pro quo tum $x$, tum $y$ valores integros recipiant,
in quo casu ex formula [1] nulli valores integri ipsorum $x$, $y$ deduci poterunt.
Si vero inter illos indices aliqui sunt, puta $\mu,\mu',\mu''$ etc. quibus valores
integri ipsorum $x$, $y$ respondent, omnes valores integri ipsorum $x$, $y$,
qui quidem ex formula [1] obtineri possunt, ii erunt, quorum indices sub aliqua
formularum $\mu+k\rho,\ \mu'+k\rho,$
\SourceLeaf{223}
$\mu''+k\rho$ etc. sunt contenti, denotante $k$ indefinite omnes numeros integros
positivos, inclusa etiam cifra.

Formulae reliquae sub quibus valores ipsorum $p$, $q$ contenti sunt, prorsus
eodem modo sunt tractandae. Si contingeret, ut ex nulla omnium harum formularum
valores integri ipsorum $x$, $y$ obtineantur, aequatio proposita in integris
nullo prorsus modo solvi posset; quoties vero revera est solubilis, omnes solutiones
in integris per praecepta in praecc. tradita exhiberi poterunt.

\Needspace{8\baselineskip}
\Article{218}
Quando $bb-ac$ est numerus quadratus atque $M=0$, omnes valores ipsorum
$p$, $q$ comprehensi erunt sub duabus huiusmodi formulis $p=\mathfrak{A}z$,
$q=\mathfrak{B}z$; $p=\mathfrak{A}'z$, $q=\mathfrak{B}'z$, ubi $z$
indefinite designat quemvis numerum integrum, $\mathfrak{A},\mathfrak{B},\mathfrak{A}',\mathfrak{B}'$
vero sunt integri dati, quorum primus cum secundo, tertius cum quarto divisorem
communem non habent (art. 212). Omnes itaque valores integri ipsorum $x$, $y$
ex formula prima oriundi contenti erunt sub formula [1]
\[
x=\frac{\mathfrak{A}z+cd-be}{bb-ac},\quad
y=\frac{\mathfrak{B}z+ae-bd}{bb-ac}
\]
omnesque reliqui ex formula secunda oriundi sub hac [2]
\[
x=\frac{\mathfrak{A}'z+cd-be}{bb-ac},\quad
y=\frac{\mathfrak{B}'z+ae-bd}{bb-ac}
\]
Sed quoniam utraque formula etiam valores fractos praebere potest (nisi $bb-ac=1$),
opus est ut eos valores ipsius $z$, qui tum ipsum $x$ tum ipsum $y$ integrum
reddunt, a reliquis in utraque formula separemus; attamen sufficit primam solam
considerare, quum pro altera prorsus eadem methodus adhibenda sit.

Quoniam $\mathfrak{A}$, $\mathfrak{B}$ inter se primi sunt, duos numeros
$\mathfrak{a}$, $\mathfrak{b}$ ita determinare licebit, ut fiat
$\mathfrak{a}\mathfrak{A}+\mathfrak{b}\mathfrak{B}=1$. Quo facto habetur
\[
(\mathfrak{a}x+\mathfrak{b}y)(bb-ac)=z+\mathfrak{a}(cd-be)+\mathfrak{b}(ae-bd)
\]
unde statim patet, omnes valores ipsius $z$ qui valores integros ipsorum $x$, $y$
producere possint, necessario numero $\mathfrak{a}(be-cd)+\mathfrak{b}(bd-ae)$
sec. mod. $bb-ac$ congruos, sive sub formula
$(bb-ac)z'+\mathfrak{a}(be-cd)+\mathfrak{b}(bd-ae)$ contentos esse debere,
designante $z'$ indefinite numerum integrum. Hinc facile loco formulae [1]
obtinemus sequentem
\SourceLeaf{224}
\[
\begin{aligned}
x&=\mathfrak{A}z'+\mathfrak{b}\times\frac{\mathfrak{A}(bd-ae)-\mathfrak{B}(be-cd)}{bb-ac}\\
y&=\mathfrak{B}z'-\mathfrak{a}\times\frac{\mathfrak{A}(bd-ae)-\mathfrak{B}(be-cd)}{bb-ac}
\end{aligned}
\]
quam aut pro omnibus valoribus ipsius $z'$ aut pro nullo valores integros ipsorum
$x,y$ praebere manifestum est, et quidem casus prior semper locum habebit, quando
$\mathfrak{A}(bd-ae)$ et $\mathfrak{B}(be-cd)$ sec. mod. $bb-ac$ sunt congrui,
posterior quando sunt incongrui. --- Prorsus eodem modo tractanda erit formula [2],
solutionesque in integris (si quas praebere potest) a reliquis separandae.

\Needspace{8\baselineskip}
\Article{219}
Quando $bb-ac=0$, forma $axx+2bxy+cyy$ exhiberi poterit ita:
$m(\alpha x+\beta y)^2$, ubi $m,\alpha,\beta$ sunt integri (art. 215).
Ponatur $\alpha x+\beta y=z$, transitque aequatio proposita in hanc:
\[
mzz+2dx+2ey+f=0
\]
unde et ex $z=\alpha x+\beta y$, deducitur
\[
x=\frac{\beta mzz+2ez+\beta f}{2\alpha e-2\beta d},\quad
y=\frac{\alpha mzz+2dz+\alpha f}{2\beta d-2\alpha e}
\]
Iam patet, nisi fuerit $\alpha e=\beta d$ (quem casum statim seorsim considerabimus),
valores ipsorum $x$, $y$, ex his formulis deductos tribuendo ipsi $z$ valorem
quemcunque, aequationi propositae satisfacere; quare nihil superest, nisi ut eos
valores ipsius $z$ determinare doceamus, ex quibus valores integri ipsorum $x$, $y$
sequantur.

Quoniam $\alpha x+\beta y=z$, necessario pro $z$ numeri \textit{integri} tantum
accipi possunt; praeterea vero manifestum est, si aliquis valor ipsius $z$ tum
ipsum $x$ tum ipsum $y$ integrum reddat, omnes valores ipsius $z$ illi secundum
modulum $2\alpha e-2\beta d$ congruos itidem valores integros producere.
Quodsi itaque pro $z$ omnes numeri integri a $0$ usque ad $2\alpha e-2\beta d-1$
(quando $\alpha e-\beta d$ est positivus) aut ad $2\beta d-2\alpha e-1$
(quando $\alpha e-\beta d$ est negativus) incl. substituuntur, et pro nullo
horum valorum tum $x$ tum $y$ integri fiunt, nullus omnino valor ipsius $z$
valores integros ipsorum $x,y$ producet, aequatioque proposita in integris nullo
modo poterit resolvi; si vero quidam ex illis valoribus ipsius $z$ ipsis $x,y$
valores integros conciliant, puta hi $\zeta,\zeta',\zeta''$ etc. (quos etiam per
solutionem congruentiarum secundi gradus ex principiis sect. IV invenire licet);
\textit{omnes} solutiones prodi%
\SourceLeaf{225}%
bunt ponendo $z=(2\alpha e-2\beta d)v+\zeta$,
$z=(2\alpha e-2\beta d)v+\zeta'$ etc., designante $v$ indefinite omnes numeros
integros.

\Needspace{8\baselineskip}
\Article{220}
Pro eo quem exclusimus casu, ubi $\alpha e=\beta d$, methodum peculiarem indagare
oportet. Supponamus, $\alpha$, $\beta$ inter se primos esse, quod licere ex
art. 215. I constat, eritque $\frac{d}{\alpha}=\frac{e}{\beta}$ numerus integer
(art. 19), quem statuemus $=h$. Tunc aequatio proposita hanc induit formam:
\[
(m\alpha x+m\beta y+h)^2-hh+mf=0
\]
manifestoque adeo rationaliter solvi nequit, nisi $hh-mf$ fuerit numerus quadratus.
Sit $hh-mf=kk$, patetque aequationi propositae sequentes duas aequivalere:
\[
m\alpha x+m\beta y+h+k=0,\quad\text{et}\quad
m\alpha x+m\beta y+h-k=0
\]
\textit{i. e.} quamlibet solutionem aequationis propositae etiam alterutri harum
aequationum satisfacere, et vice versa. Aequatio prior manifesto in integris solvi
nequit, nisi $h+k$ per $m$ fuerit divisibilis, similiterque posterior solutionem
in integris non admittet, nisi $h-k$ per $m$ fuerit divisibilis. Hae vero
conditiones ad resolubilitatem utriusque aequationis sufficiunt (quia $\alpha$,
$\beta$ inter se primi esse supponuntur), omnesque solutiones secundum regulas
notas exhiberi poterunt.

\Needspace{8\baselineskip}
\Article{221}
Casum in art. 217 consideratum (quia omnium difficillimus est) exemplo illustramus.
Proposita sit aequatio $xx+8xy+yy+2x-4y+1=0$. Ex hac primo per introductionem
aliarum incognitarum $p=15x-9$. $q=15y+6$ derivatur aequatio
$pp+8pq+qq=-540$. Huius autem solutiones omnes in integris, contineri inveniuntur
sub quatuor formulis sequentibus:
\[
\begin{aligned}
p&=6t,&q&=-24t-90u\\
p&=6t,&q&=-24t+90u\\
p&=-6t,&q&=24t-90u\\
p&=-6t,&q&=24t+90u
\end{aligned}
\]
denotantibus $t$, $u$ indefinite omnes numeros integros positivos aequationi
\SourceLeaf{226}
$tt-15uu=1$ satisfacientes, quos complectitur formula:
\[
\begin{aligned}
t&=\tfrac{1}{2}\bigl((4+\sqrt{15})^n+(4-\sqrt{15})^n\bigr)\\
u&=\frac{1}{2\sqrt{15}}\bigl((4+\sqrt{15})^n-(4-\sqrt{15})^n\bigr)
\end{aligned}
\]
si $n$ indefinite omnes numeros integros positivos (inclusa etiam cifra) designat.
Quamobrem omnes valores ipsorum $x,y$ contenti erunt sub formulis his:
\[
\begin{aligned}
x&=\tfrac{1}{5}(2t+3),&y&=-\tfrac{1}{5}(8t+30u+2)\\
x&=\tfrac{1}{5}(2t+3),&y&=-\tfrac{1}{5}(8t-30u+2)\\
x&=\tfrac{1}{5}(-2t+3),&y&=\tfrac{1}{5}(8t-30u-2)\\
x&=\tfrac{1}{5}(-2t+3),&y&=\tfrac{1}{5}(8t+30u-2)
\end{aligned}
\]
Praeceptis autem nostris rite applicatis, reperietur, ut valores \textit{integri}
prodeant, in formula prima et secunda eos valores ipsorum $t,u$ accipi debere,
qui proveniant ex indice $n$ \textit{pari}; in tertia quartaque vero eos,
qui ex \textit{impari} $n$ obtineantur. --- Solutiones simplicissimae habentur
hae: $x=1,-1,-1$; $y=-2,0,12$ resp.

Ceterum observare convenit, solutionem problematis in artt. praecc. explicati
plerumque per multifaria artificia abbreviari posse, praesertim quantum ad
exclusionem solutionum inutilium \textit{i. e.} fractiones implicantium pertinet;
sed haec ne nimis longi fiamus hoc loco praeterire coacti sumus.

\Needspace{12\baselineskip}
\begin{center}
\textit{Annotationes historicae.}
\end{center}

\Needspace{8\baselineskip}
\Article{222}
Quoniam complura ex iis quae hucusque pertractavimus etiam ab aliis geometris
considerata sunt, horum merita silentio praeterire non possumus. De \textit{formarum
aequivalentia} disquisitiones generales instituit ill. La Grange, \textit{Nouv.
M\'em. de l'Ac. de Berlin,} 1773 \textit{p.} 263 et 1775 \textit{p.} 323 sqq.,
ubi imprimis docuit, pro quovis determinante dato multitudinem finitam formarum
dari ita comparatarum, ut quaevis forma illius determinantis alicui ex ipsis
aequivalens sit, adeoque omnes formas determinantis dati in classes distribui posse.
Postea clar. Le Gendre plures proprietates elegantes huius classificationis ad
maximam partem per inductionem detexit, quas infra trademus demonstrationibusque
muniemus. Ceterum distinctionem aequivalentiae propriae et impropriae, cuius usus
maxime in disquisitionibus subtilioribus conspicuus est, nemo hucusque attigerat.

Problema famosum in art. 216 sqq. explicatum ill. La Grange primus com%
\SourceLeaf{227}%
plete resolvit, \textit{Hist. de l'Ac. de Berlin,} 1767 \textit{p.} 165 et
1768 \textit{p.} 181 sqq. Exstat solutio (sed minus completa) etiam in
\textit{Suppl. ad Euleri Algebram} iam saepius laudatis. Iam antea ill. Euler
idem argumentum aggressus fuerat, \textit{Comm. Petr. T.} VI \textit{p.} 175;
\textit{Comm. Nov. T.} IX \textit{p.} 3; \textit{Ibid. T.} XVIII \textit{p.} 185
sqq., sed investigationem suam eo semper restrinxit, ut ex aliqua solutione,
quam iam cognitam esse supponit, aliae deriventur; praetereaque ipsius methodi in
paucis tantummodo casibus \textit{omnes} solutiones suppeditare valent
(vid. La Grange \textit{Hist. de l'Ac. de Berlin} 1767, \textit{p.} 237).
Quum ultima harum trium commentt. recentioris dati sit quam solutio La Grangiana,
quae problema omni generalitate amplectitur nihilque hoc respectu desiderandum
relinquit: Euler tunc temporis (Tomus XVIII Commentariorum pertinet ad annum 1773,
et a. 1774 est publicatus) illam solutionem nondum novisse videtur. Ceterum solutio
nostra (perinde ut omnia reliqua quae in hac sectione hactenus tradidimus),
principiis omnino diversis est superstructa.

Quae ab aliis, Diophanto, Fermatio etc. huc pertinentia sunt tradita, casus maxime
speciales spectant; quare quum eorum quae praesertim memoratu digna visa sunt,
iam supra mentio facta sit, sigillatim omnia enarrare supersedemus.

\begin{center}
$\begin{array}{ccc}\ast&&\ast\\&\ast&\end{array}$
\end{center}

Quae hactenus de formis secundi gradus exposuimus, pro primis tantum elementis
huius doctrinae sunt habenda: innixius hanc disquisitionem persequentibus campus
se aperuit nobis vastissimus, ex quo ea quae attentione imprimis digna videntur,
in sequentibus excerpemus. Namque argumentum hoc tam fertile est, ut permulta alia,
quae iam nunc invenire nobis contigit, brevitatis gratia silentio praeterire
oporteat: multo vero plura sine dubio adhuc latent novosque conatus exspectant.
Ceterum in limine harum investigationum statim adnotare convenit, formas
determinantis $0$ inde exclusas esse, nisi contrarium moneatur.

\Needspace{12\baselineskip}
\begin{center}
\textit{DISQUISITIONES ULTERIORES DE FORMIS.}
\end{center}

\Needspace{12\baselineskip}
\begin{center}
\textit{Distributio formarum determinantis dati in classes.}
\end{center}

\Needspace{8\baselineskip}
\Article{223}
Iam supra (artt. 175, 195, 211) ostendimus, proposito numero quocunque integro
$D$ (sive positivo sive negativo) assignari posse multitudinem finitam formarum
$F,F',F''$ etc. determinantis $D$, ita comparatarum, ut quaevis forma de%
\SourceLeaf{228}%
terminantis $D$ proprie aequivalens sit alicui ex illis et quidem unicae tantum.
Omnes igitur formae determinantis $D$ (quarum multitudo est infinita) secundum
illas formas \textit{classificari} poterunt, formando scilicet e complexu omnium
formarum formae $F$ proprie aequivalentium classem primam; e formis quae formae
$F'$ proprie aequivalent, secundam etc.

Ex singulis classibus formarum determinantis dati $D$, forma aliqua eligi et
tamquam \textit{forma repraesentans} totius classis considerari poterit. Per se
quidem prorsus arbitrarium est, quaenam forma ex quaque classe accipiatur, attamen
ea semper praeferenda erit, quae reliquas \textit{simplicitate} superare videtur.
Simplicitas formae alicuius $(a,b,c)$ manifesto ex magnitudine numerorum
$a,b,c$ aestimanda est, meritoque forma $(a',b',c')$ minus simplex dicetur
quam $(a,b,c)$ si $a'>a$, $b'>b$, $c'>c$. Sed hinc res nondum
determinatur penitus, arbitrioque nostro relinquitur \textit{e. g.}, utram ex
formis $(17,0,-45)$, $(5,0,-153)$ pro simpliciori habere malimus.
Plerumque tamen e re erit, sequentem normam observare:

I. Quando determinans $D$ est negativus, adoptentur formae reductae in singulis
classibus contentae tamquam formae repraesentantes; ubi vero in eadem classe duae
formae reductae reperiuntur (quae erunt oppositae, art. 172). recipiatur ea,
cuius terminus medius positivus.

II. Quando determinans $D$ est positivus non-quadratus, evolvatur periodus formae
alicuius reductae in classe proposita contentae, in qua aut duae formae ancipites
invenientur aut nulla (art. 187).

1) In casu priori sint formae ancipites hae: $(A,B,C)$, $(A',B',C')$;
residua minima numerorum $B,B'$ secundum modulos $A,A'$ resp. $M,M'$
(quae positive accipi poterunt nisi sunt $=0$); denique
$\frac{D-MM}{A}=N$, $\frac{D-M'M'}{A'}=N'$.
His ita factis, ex formis $(A,M,-N)$, $(A',M',-N')$ ea quae simplicissima
videtur, pro forma repraesentante accipiatur. In hoc iudicio forma cuius terminus
medius $=0$, praeferatur; quando vero terminus medius aut in utraque aut in
neutra est $0$, ea quae terminum primum minorem habet, alteri praehabenda,
et quando termini primi magnitudine sunt aequales signis diversi, signum negativum
positivo postponendum.

2) Quando vero nulla forma anceps in tota periodo habetur. eligatur ex omnibus
periodi formis ea quae terminum primum sine respectu signi minimum
\SourceLeaf{229}
habet, ita quidem, ut si duae formae in eadem periodo occurrant, in quarum altera
idem terminus primus signo positivo affectus sit in altero negativo, posterior
priori postponatur. Sit haec forma $(A,B,C)$, deducaturque ex ipsa eodem modo
ut in casu praec. forma alia $(A,M,-N)$ (puta, accipiendo pro $M$ residuum
absolute minimum ipsius $B$ secundum mod. $A$, et faciendo
$N=\frac{D-MM}{A}$): haec demum pro repraesentante adoptetur.

Quodsi vero eveniret, ut idem terminus primus minimus $A$ pluribus periodi formis
communis sit, omnes hae formae eo quo praescripsimus modo tractandae et ex formis
prodeuntibus ea cuius terminus medius quam minimus evadit tamquam forma
repraesentans assumenda erit.

Ita \textit{e. g.} pro $D=305$ habetur periodus inter alias haec:
$(17,4,-17)$, $(-17,13,8)$, $(8,11,-23)$, $(-23,12,7)$,
$(7,16,-7)$, $(-7,12,23)$, $(23,11,-8)$, $(-8,13,17)$,
ex qua primo eligitur forma $(7,16,-7)$, hincque secundo deducitur forma
repraesentans $(7,2,-43)$.

III. Quando determinans est positivus quadratus $=kk$, eruatur forma reducta
$(A,k,0)$ in classe proposita contenta et, si $A<k$ aut $=k$, pro forma
repraesentante ipsa recipiatur; si vero $A>k$, assumatur illius loco forma
$(A-2k,k,0)$, cuius terminus primus erit negativus, sed minor quam $k$.

\textit{Ex.} Hoc modo omnes formae determinantis $-235$ distribuuntur in classes
sedecim, quarum repraesentantes erunt: $(1,0,235)$, $(2,1,118)$,
$(4,1,59)$, $(4,-1,59)$, $(5,0,47)$, $(10,5,26)$,
$(13,5,20)$, $(13,-5,20)$, octoque aliae a praecedentibus in solis signis
terminorum externorum diversae $(-1,0,-235)$, $(-2,1,-118)$ etc.

Omnes formae determinantis $79$ in sex classes discedunt, quarum repraesentantes
$(1,0,-79)$, $(3.\,1,-26)$, $(3,-1,-26)$, $(-1,0.\,79)$,
$(-3,1,26)$, $-3.\,-1,26)$.

\Needspace{8\baselineskip}
\Article{224}
Per hanc itaque classificationem formae quae proprie aequivalentes sunt. a reliquis
omnino segregabuntur. Duae formae eiusdem determinantis $D$, si ex eadem classe
sunt, proprie aequivalentes erunt; quivis numerus per unam repraesentabilis etiam
per alteram repraesentari poterit; et si numerus quicunque $M$ per formam priorem
ita repraesentari potest. ut indeterminatae valores inter se
\SourceLeaf{230}
primos habeant, idem numerus per alteram formam eodem modo repraesentari poterit,
et quidem ita, ut utraque repraesentatio ad eundem valorem expressionis
$\sqrt{D}\ (\mathrm{mod.}\ M)$ pertineat. Si vero duae formae ad classes diversas
pertinent, proprie aequivalentes non erunt; a repraesentabilitate numeri alicuius
dati per unam ad repraesentabilitatem eiusdem numeri per alteram concludi nequit;
contra, si numerus $M$ per alteram repraesentari potest ita ut valores
indeterminatarum inter se primi sint, statim certi sumus, nullam similem
repraesentationem eiusdem numeri per formam alteram dari, quae ad eundem valorem
expr. $\sqrt{D}(\mathrm{mod.}\ M)$ pertineat (V. artt. 167, 168).

Contra utique fieri potest, ut formae duae $F,F'$, e classibus diversis $K,K'$
improprie aequivalentes sint, in quo casu \textit{quaevis} forma ex altera classe
\textit{cuivis} formae ex altera improprie aequivalebit; quaevis forma ex $K$
formam sibi oppositam habebit in $K'$, classesque ipsae $K,K'$ \textit{oppositae}
dicentur. Ita in exemplo primo art. praec. classis tertia formarum det. $-235$
quartae, septima octavae opposita est; in ex. secundo classis secunda tertiae,
quinta sextae. Propositis itaque duabus formis quibuscunque e classibus oppositis,
quivis numerus $M$ qui per alteram repraesentari potest, etiam per alteram
poterit; quod, si in altera fit per valores indeterminatarum inter se primos,
in altera perinde fieri poterit, ita tamen, ut hae duae repraesentationes ad valores
oppositos expr. $\sqrt{D}\ (\mathrm{mod.}\ M)$ pertineant. --- Ceterum regulae
supra traditae pro electione formarum repraesentantium ita sunt constitutae,
ut classes oppositae formas repraesentantes oppositas semper nanciscantur.

Denique dantur etiam classes \textit{sibi ipsis oppositae.} Scilicet si forma aliqua
simul cum forma opposita in eadem classe continetur, facile perspicitur, omnes
formas huius classis tum proprie tum improprie inter se aequivalentes esse,
oppositasque suas secum habere. Hanc indolem quaevis classis habebit, in qua forma
anceps continetur, et vice versa in quavis classe sibi ipsi opposita necessario
forma anceps reperietur (art. 163, 165), quamobrem \textit{classis anceps}
nuncupabitur. Ita inter classes formarum determinantis $-235$ octo ancipites
habentur, quarum repraesentantes sunt $(1,0,235)$, $(2,1,118)$, $(5,0,47)$,
$(10,5,26)$, $(-1,0,-235)$, $(-2,1,-118)$, $(-5,0,-47)$,
$(-10,5,-26)$; inter classes formarum determinantis $79$ duae, quarum
repraesentantes $(1,0,-79)$, $(-1,0,79)$. --- Ceterum si formae
repraesentantes secundum regulas nostras determinatae sunt, classes ancipites
nullo negotio inde cognosci poterunt. Scilicet pro determinante positivo
\SourceLeaf{231}
non-quadrato classis anceps certo formam repraesentantem ancipitem nanciscitur
(art. 194); pro determinante negativo forma repraesentans classis ancipitis aut
ipsa anceps erit, aut talis cuius termini externi sunt aequales (art. 172);
denique pro determinante positivo quadrato per art. 210 facile diiudicatur,
an forma repraesentans sibi ipsi improprie aequivalens sit adeoque classis,
quam repraesentat, anceps.

\Needspace{8\baselineskip}
\Article{225}
Iam supra (art. 175) ostendimus, in forma $(a,b,c)$ determinantis negativi
terminos externos eadem signa habere tum inter se tum cum terminis externis
cuiusvis aliae formae illi aequivalentis. Si $a,c$ sunt positivi, formam
$(a,b,c)$ \textit{positivam} vocabimus, nec non totam classem in qua $(a,b,c)$
continetur et quae e solis formis positivis constabit, \textit{classem positivam}
dicemus. Contra $(a,b,c)$ erit \textit{forma negativa,} et in \textit{classe
negativa} contenta, si $a,c$ sunt negativi. Per formam positivam numeri negativi,
per negativam positivi repraesentari nequeunt. Si forma $(a,b,c)$ est
repraesentans alicuius classis positivae, forma $(-a,b,-c)$ repraesentans classis
negativae erit, unde sequitur, multitudinem classium positivarum multitudini
negativarum aequalem esse, et, simul ac illae fuerint assignatae, etiam has haberi.
Quocirca in disquisitionibus super formis determinantis negativi plerumque sufficit
classes positivas considerare, quippe quarum proprietates ad classes negativas
facile transferuntur.

Ceterum distinctio haec unice in formis determinantis negativi locum habet;
per formas determinantis positivi sine discrimine numeri positivi et negativi
repraesentari possunt, quin adeo haud raro duae formae tales $(a,b,c)$,
$(-a,b,-c)$ in hoc casu ad eandem classem sunt referendae.

\Needspace{12\baselineskip}
\begin{center}
\textit{Distributio classium in ordines.}
\end{center}

\Needspace{8\baselineskip}
\Article{226}
Formam quamcunque $(a,b,c)$ \textit{primitivam} vocamus, si numeri $a,b,c$
divisorem communem non habent; alioquin dicitur \textit{derivata,} et quidem,
posito numerorum $a,b,c$ divisore communi maximo $=m$, forma $(a,b,c)$
erit \textit{derivata e forma primitiva} $\left(\frac{a}{m},\frac{b}{m},\frac{c}{m}\right)$.
Ex hac definitione statim liquet, omnes formas, quarum determinans per nullum
quadratum (praeter $1$) divisibilis sit, necessario primitivas esse. Porro ex
art. 161 patet, si in aliqua classe data formarum determinantis $D$ forma
primitiva inveniatur, omnes formas huius classis primitivas fore, in quo
